%%%%%%%%%%%%%%%%%%%%%%%%%%%%%%%%%%%%%%%%%%%%%%%%%%%%%%%%%%%%%%%%%%%%%%%%
%                                                                      %
%     File: Thesis_Abstract.tex                                        %
%     Tex Master: Thesis.tex                                           %
%                                                                      %
%     Author: Andre C. Marta                                           %
%     Last modified :  4 Mar 2024                                      %
%                                                                      %
%%%%%%%%%%%%%%%%%%%%%%%%%%%%%%%%%%%%%%%%%%%%%%%%%%%%%%%%%%%%%%%%%%%%%%%%

\section*{Abstract}

% Add entry in the table of contents as section
\addcontentsline{toc}{section}{Abstract}

An unmanned surface vessel (USV) and an unmanned aerial vehicle (UAV) connected by a tether combine the endurance of the USV with the view from above of the UAV, which draws power and a reliable communication link through the tether. The tether, however, couples the two vehicles, which must remain within its reach while each respects its own limits, under waves and wind. This thesis develops the cooperative control of such a system, formed by the WAM-V USV and the x500 UAV. Models of both vehicles are derived and validated against a simulator, together with a closed-form model of the force of a slack tether that accounts for its weight. Each vehicle is driven by its own non-linear model predictive controller, and the two controllers exchange their predicted trajectories, so that each anticipates the motion of the other within the reach of the tether. A disturbance observer compensates unmodelled forces, such as the wind. In a simulation at sea, the UAV controller tracks a trajectory with less than a quarter of the error of the PX4 position controller, and both controllers keep their vehicles within their limits when the reference cannot be followed. The UAV lands on the moving deck a few centimetres from the centre of the pad in waves of up to $1.2\,\text{m}$ and winds of up to $8\,\text{m/s}$, and in missions of tower inspection, target pursuit and UAV recovery, the two vehicles operate as a single system.

\vfill

\textbf{\Large Keywords:} Tether, Unmanned Aerial Vehicle, Unmanned Surface Vessel, Cooperative Control, Non-linear Model Predictive Control

